% Options for packages loaded elsewhere
\PassOptionsToPackage{unicode}{hyperref}
\PassOptionsToPackage{hyphens}{url}
\PassOptionsToPackage{dvipsnames,svgnames,x11names}{xcolor}
%
\documentclass[
  10pt,
]{article}
\usepackage{amsmath,amssymb}
\usepackage{iftex}
\ifPDFTeX
  \usepackage[T1]{fontenc}
  \usepackage[utf8]{inputenc}
  \usepackage{textcomp} % provide euro and other symbols
\else % if luatex or xetex
  \usepackage{unicode-math} % this also loads fontspec
  \defaultfontfeatures{Scale=MatchLowercase}
  \defaultfontfeatures[\rmfamily]{Ligatures=TeX,Scale=1}
\fi
\usepackage{lmodern}
\ifPDFTeX\else
  % xetex/luatex font selection
  \setmainfont[]{DejaVu Serif}
  \setsansfont[]{DejaVu Sans}
  \setmonofont[]{DejaVu Sans Mono}
  \setmathfont[]{Latin Modern Math}
\fi
% Use upquote if available, for straight quotes in verbatim environments
\IfFileExists{upquote.sty}{\usepackage{upquote}}{}
\IfFileExists{microtype.sty}{% use microtype if available
  \usepackage[]{microtype}
  \UseMicrotypeSet[protrusion]{basicmath} % disable protrusion for tt fonts
}{}
\makeatletter
\@ifundefined{KOMAClassName}{% if non-KOMA class
  \IfFileExists{parskip.sty}{%
    \usepackage{parskip}
  }{% else
    \setlength{\parindent}{0pt}
    \setlength{\parskip}{6pt plus 2pt minus 1pt}}
}{% if KOMA class
  \KOMAoptions{parskip=half}}
\makeatother
\usepackage{xcolor}
\usepackage[margin=2.0cm]{geometry}
\usepackage{color}
\usepackage{fancyvrb}
\newcommand{\VerbBar}{|}
\newcommand{\VERB}{\Verb[commandchars=\\\{\}]}
\DefineVerbatimEnvironment{Highlighting}{Verbatim}{commandchars=\\\{\}}
% Add ',fontsize=\small' for more characters per line
\newenvironment{Shaded}{}{}
\newcommand{\AlertTok}[1]{\textcolor[rgb]{1.00,0.00,0.00}{\textbf{#1}}}
\newcommand{\AnnotationTok}[1]{\textcolor[rgb]{0.38,0.63,0.69}{\textbf{\textit{#1}}}}
\newcommand{\AttributeTok}[1]{\textcolor[rgb]{0.49,0.56,0.16}{#1}}
\newcommand{\BaseNTok}[1]{\textcolor[rgb]{0.25,0.63,0.44}{#1}}
\newcommand{\BuiltInTok}[1]{\textcolor[rgb]{0.00,0.50,0.00}{#1}}
\newcommand{\CharTok}[1]{\textcolor[rgb]{0.25,0.44,0.63}{#1}}
\newcommand{\CommentTok}[1]{\textcolor[rgb]{0.38,0.63,0.69}{\textit{#1}}}
\newcommand{\CommentVarTok}[1]{\textcolor[rgb]{0.38,0.63,0.69}{\textbf{\textit{#1}}}}
\newcommand{\ConstantTok}[1]{\textcolor[rgb]{0.53,0.00,0.00}{#1}}
\newcommand{\ControlFlowTok}[1]{\textcolor[rgb]{0.00,0.44,0.13}{\textbf{#1}}}
\newcommand{\DataTypeTok}[1]{\textcolor[rgb]{0.56,0.13,0.00}{#1}}
\newcommand{\DecValTok}[1]{\textcolor[rgb]{0.25,0.63,0.44}{#1}}
\newcommand{\DocumentationTok}[1]{\textcolor[rgb]{0.73,0.13,0.13}{\textit{#1}}}
\newcommand{\ErrorTok}[1]{\textcolor[rgb]{1.00,0.00,0.00}{\textbf{#1}}}
\newcommand{\ExtensionTok}[1]{#1}
\newcommand{\FloatTok}[1]{\textcolor[rgb]{0.25,0.63,0.44}{#1}}
\newcommand{\FunctionTok}[1]{\textcolor[rgb]{0.02,0.16,0.49}{#1}}
\newcommand{\ImportTok}[1]{\textcolor[rgb]{0.00,0.50,0.00}{\textbf{#1}}}
\newcommand{\InformationTok}[1]{\textcolor[rgb]{0.38,0.63,0.69}{\textbf{\textit{#1}}}}
\newcommand{\KeywordTok}[1]{\textcolor[rgb]{0.00,0.44,0.13}{\textbf{#1}}}
\newcommand{\NormalTok}[1]{#1}
\newcommand{\OperatorTok}[1]{\textcolor[rgb]{0.40,0.40,0.40}{#1}}
\newcommand{\OtherTok}[1]{\textcolor[rgb]{0.00,0.44,0.13}{#1}}
\newcommand{\PreprocessorTok}[1]{\textcolor[rgb]{0.74,0.48,0.00}{#1}}
\newcommand{\RegionMarkerTok}[1]{#1}
\newcommand{\SpecialCharTok}[1]{\textcolor[rgb]{0.25,0.44,0.63}{#1}}
\newcommand{\SpecialStringTok}[1]{\textcolor[rgb]{0.73,0.40,0.53}{#1}}
\newcommand{\StringTok}[1]{\textcolor[rgb]{0.25,0.44,0.63}{#1}}
\newcommand{\VariableTok}[1]{\textcolor[rgb]{0.10,0.09,0.49}{#1}}
\newcommand{\VerbatimStringTok}[1]{\textcolor[rgb]{0.25,0.44,0.63}{#1}}
\newcommand{\WarningTok}[1]{\textcolor[rgb]{0.38,0.63,0.69}{\textbf{\textit{#1}}}}
\usepackage{longtable,booktabs,array}
\usepackage{calc} % for calculating minipage widths
% Correct order of tables after \paragraph or \subparagraph
\usepackage{etoolbox}
\makeatletter
\patchcmd\longtable{\par}{\if@noskipsec\mbox{}\fi\par}{}{}
\makeatother
% Allow footnotes in longtable head/foot
\IfFileExists{footnotehyper.sty}{\usepackage{footnotehyper}}{\usepackage{footnote}}
\makesavenoteenv{longtable}
\setlength{\emergencystretch}{3em} % prevent overfull lines
\providecommand{\tightlist}{%
  \setlength{\itemsep}{0pt}\setlength{\parskip}{0pt}}
\setcounter{secnumdepth}{-\maxdimen} % remove section numbering
% doc_header.tex -- LaTeX preamble additions used by build_docs.sh for every
% Markdown -> LaTeX -> PDF conversion (the documents of the folder Revision_SM_Higgs).
\usepackage{etoolbox}
% longtable may not be loaded by pandoc when a document has no table
\usepackage{longtable}
% bold vectors: \bm -> \symbf under unicode-math (XeLaTeX), bm package otherwise
\makeatletter
\@ifpackageloaded{unicode-math}{\providecommand{\bm}[1]{\symbf{#1}}}{\usepackage{bm}}
\makeatother
% smaller type in long tables so that wide audit tables fit the text width;
% inside a table a line break is allowed after an underscore of plain text
% (pandoc writes it as \_), so that long identifiers wrap inside their column
% rows are set a little apart, so that symbols with indices do not touch
\AtBeginEnvironment{longtable}{\footnotesize\renewcommand{\arraystretch}{1.45}%
  \let\sptunderscore\_\def\_{\sptunderscore\allowbreak}}
\setlength{\LTcapwidth}{\textwidth}
% A short table is kept on one page.  The build (code/scripts/keep_tables.py) ends the rows of such a table with \\*,
% which forbids a page break after a row, and sets \sptkeptpenalty to 10000 in front of the table.  This penalty is
% put here in the place of the breakpoint (penalty 0) that longtable leaves after every table, so that the paragraph
% that follows a kept table, its caption, begins on the same page; it is set back to 0 after the table.
\newcount\sptkeptpenalty
\sptkeptpenalty=0
\makeatletter
\patchcmd{\endlongtable}{\endgraf\penalty\z@\addvspace\LTpost}%
  {\endgraf\penalty\sptkeptpenalty\addvspace\LTpost\global\sptkeptpenalty=\z@}%
  {}{\PackageError{doc_header}{the end of longtable could not be patched}{The tables can still be typeset, but a short table may be separated from its caption.}}
\makeatother
% allow line breaks after slashes and hyphens inside table cells
\usepackage{xurl}
% code blocks: break over-long lines instead of running into the margin.
% Plain ``` blocks are typeset by pandoc as {verbatim}, highlighted ones
% (```bash) as {Highlighting}; both are re-declared as breakable fancyvrb
% environments.  The continuation line is marked with a small hooked arrow.
\usepackage{fvextra}
\RecustomVerbatimEnvironment{verbatim}{Verbatim}{breaklines,breakanywhere,fontsize=\small}
\makeatletter
\@ifundefined{Highlighting}{}{%
  \RecustomVerbatimEnvironment{Highlighting}{Verbatim}{breaklines,breakanywhere,fontsize=\small,commandchars=\\\{\}}}
\makeatother
% inline code (pandoc writes it as \texttt{...}): never hyphenate -- an inserted
% hyphen could be read as part of a file name -- but allow a line break after
% an underscore, slash, comma, colon or equals sign, so that long paths and
% column lists wrap inside the text width and inside table columns.
% Inside a formula (e.g. the name of a notebook symbol in an equation) the same
% command must first leave math mode, which \text of amsmath does.
\ExplSyntaxOn
\tl_new:N \l__spt_code_tl
\cs_new_protected:Npn \__spt_texttt:n #1
  {
    \tl_set:Nn \l__spt_code_tl {#1}
    \regex_replace_all:nnN { ([\/\,\:\=]) } { \1 \c{allowbreak} } \l__spt_code_tl
    \regex_replace_all:nnN { \c{_} } { \c{_} \c{allowbreak} } \l__spt_code_tl
    \group_begin: \ttfamily \hyphenchar \font = -1 \scan_stop: \l__spt_code_tl \group_end:
  }
\RenewDocumentCommand \texttt { m }
  {
    \mode_if_math:TF { \text { \__spt_texttt:n {#1} } } { \__spt_texttt:n {#1} }
  }
\ExplSyntaxOff
% A figure is written in the Markdown source as an image paragraph followed by
% a caption paragraph ("**Figure S1.** ..."), so that the caption is also
% visible on GitHub.  An image that starts a paragraph is centered.
\usepackage{graphicx}% pandoc loads it only for documents that contain an image
\NewCommandCopy\sptincludegraphics\includegraphics
\RenewDocumentCommand\includegraphics{O{} m}{%
  \ifvmode\begin{center}\sptincludegraphics[#1]{#2}\end{center}%
  \else\sptincludegraphics[#1]{#2}\fi}
% keep figures where they are referenced
\usepackage{float}
\floatplacement{figure}{H}
% running header with the document title
\usepackage{fancyhdr}
\pagestyle{fancy}
\fancyhf{}
\fancyhead[L]{\small PROJECT spacetimes, Revision\_SM\_Higgs}
\fancyhead[R]{\small\thepage}
\renewcommand{\headrulewidth}{0.4pt}
% 16x16 matrices: tighter columns so that a full matrix fits the text width
\setlength{\arraycolsep}{3.2pt}
% The documents of this project contain long tagged symbols, file paths and
% identifiers that cannot be hyphenated.  A paragraph is therefore allowed to
% have looser lines before TeX gives up and lets a line stick out of the margin.
\AtBeginDocument{\tolerance=2000 \emergencystretch=5em \hbadness=2000}
\ifLuaTeX
  \usepackage{selnolig}  % disable illegal ligatures
\fi
\IfFileExists{bookmark.sty}{\usepackage{bookmark}}{\usepackage{hyperref}}
\IfFileExists{xurl.sty}{\usepackage{xurl}}{} % add URL line breaks if available
\urlstyle{same}
\hypersetup{
  pdftitle={paper\_prd: the Regular Article ``Higgs as a Composite, and Quantum Gravity'' for Physical Review D},
  pdfauthor={Prepared for Patrick Lee Nash, Ph.D.},
  colorlinks=true,
  linkcolor={blue},
  filecolor={Maroon},
  citecolor={blue},
  urlcolor={blue},
  pdfcreator={LaTeX via pandoc}}

\title{paper\_prd: the Regular Article ``Higgs as a Composite, and
Quantum Gravity'' for Physical Review D}
\usepackage{etoolbox}
\makeatletter
\providecommand{\subtitle}[1]{% add subtitle to \maketitle
  \apptocmd{\@title}{\par {\large #1 \par}}{}{}
}
\makeatother
\subtitle{README of the folder Revision\_SM\_Higgs/paper\_prd of PROJECT
spacetimes}
\author{Prepared for Patrick Lee Nash, Ph.D.}
\date{November 22, 2026}

\begin{document}
\maketitle

\hypertarget{what-this-folder-is}{%
\section{1. What this folder is}\label{what-this-folder-is}}

This folder contains the Regular Article ``Higgs as a Composite, and
Quantum Gravity'' by Patrick Lee Nash, Ph.D., prepared for submission to
Physical Review D, together with its Supplemental Material, the cover
letter, the information to be entered on the submission server of the
American Physical Society (APS), and the folder of the files to be
uploaded. The steps on the server (logging in, entering the data,
uploading, submitting) are the author's; nothing has been submitted.

The article is the third part of stage 6 of the revision of the
hypotheses of PROJECT spacetimes. The author asked on 21 November 2026,
at 17:09 PST, that the hypotheses be modified so that the superfluid can
be the Higgs field of the Standard Model, with the same quantum numbers
as the Higgs; at 19:38 PST he instructed that a second singlet with the
hypercharge \(-\tfrac34\) be added. The article reports the second
version of the revised hypotheses, with the first version (one singlet)
as the comparison. It replaces the article that this folder held under
the earlier hypotheses, in which the spinor field had no electroweak
quantum numbers and the superfluid was the Higgs field of a new gauge
group; that article stays unchanged in the folder
\texttt{Revision\_HS\_CJT\_and\_Kohn-Sham/paper\_prd} and in the history
of the repository.

The article answers, under the revised hypotheses, the questions of the
author's request of 21 November 2026, 11:29 PST: whether the superfluid
can play the role of the Higgs particle in each of the five mean-field
methods, and whether the Higgs is then a composite particle. The first
question of that request, the width of the region of the fluctuations
near the critical temperature, was answered for the earlier hypotheses
only. The revision has no stage at finite temperature, so the article
does not recompute it and says so (decision D35 of the task record; the
author can ask for a thermal stage). The gravitational sector of the
title is not computed: the article outlines, in its Sec. X, how it will
be added through E. A. Lord's gauge theory of the local change of spinor
representation (Proc. Cambridge Philos. Soc. 69, 423 (1971)), the route
that the author named, and says that this sector is a program.

\hypertarget{the-answers-in-brief}{%
\section{2. The answers in brief}\label{the-answers-in-brief}}

The symbols are those of the article: \(\Lambda_{\rm LL}\), the momentum
cutoff of the uniform states, is the unit of energy; \(v_{\rm SH}\) is
the scale of the breaking, defined by \(M_W=gv_{\rm SH}/2\) as the
Standard Model defines \(v\); \(M_H\) is the mass of the amplitude mode.

\begin{enumerate}
\def\labelenumi{\arabic{enumi}.}
\tightlist
\item
  \textbf{The quantum numbers and the gauge bosons.} The two pair fields
  of the isodoublet with the two singlets form a bidoublet of
  \(SU(2)_L\times SU(2)_{\rm R}\) with the quantum numbers of two Higgs
  doublets. In all five methods the ground state is the custodial
  vacuum: it breaks \(SU(2)_L\times U(1)_Y\) to the \(U(1)\) of the
  electric charge, keeps the photon massless and gives \(\rho=1\). With
  \(v_{\rm SH}=246.22\) GeV the \(W\) and \(Z\) bosons have their
  measured masses at tree level and the cutoff is \(2207\) to \(2413\)
  GeV.
\item
  \textbf{The Higgs particle.} The amplitude mode is a composite scalar,
  bound below the energy that breaks a pair in the methods without
  exchange self-energies and decaying above it in KS and CJT. In the
  pairing-only methods it lies exactly at the pair self-energy:
  \(M_H/v_{\rm SH}=7.967\), against \(m_H/v=0.5082\), and no coupling of
  the scan brings it below \(4.892\). Within the mean field the
  superfluid is therefore the Higgs field of the Standard Model in its
  quantum numbers and its gauge-boson masses, but not in the mass of its
  boson.
\item
  \textbf{Compositeness.} Yes, by construction: the order parameter is
  the expectation value of pairs of a member of the isodoublet and a
  singlet, and the amplitude mode is a bound state or a resonance of two
  quasi-particles.
\item
  \textbf{What comes with it.} A custodial triplet of massive scalars, a
  massless Goldstone boson of the particle number (decision D32 of the
  author, pending), and fermions with the electric charges \(\tfrac34\)
  and \(-\tfrac14\).
\end{enumerate}

\hypertarget{contents}{%
\section{3. Contents}\label{contents}}

\begin{longtable}[]{@{}
  >{\raggedright\arraybackslash}p{(\columnwidth - 2\tabcolsep) * \real{0.5000}}
  >{\raggedright\arraybackslash}p{(\columnwidth - 2\tabcolsep) * \real{0.5000}}@{}}
\toprule\noalign{}
\begin{minipage}[b]{\linewidth}\raggedright
file or folder
\end{minipage} & \begin{minipage}[b]{\linewidth}\raggedright
what it is
\end{minipage} \\
\midrule\noalign{}
\endhead
\bottomrule\noalign{}
\endlastfoot
\texttt{higgs\_composite\_qg.pdf}, \texttt{.tex}, \texttt{.md} & the
article (REVTeX 4.2, two columns), its LaTeX source and its Markdown
version \\
\texttt{higgs\_composite\_qg\_SM.pdf}, \texttt{.tex}, \texttt{.md} & the
Supplemental Material: notation, the model and its group theory, the
methods, the normal vacuum, the stationary states, the effective
potential, the susceptibilities and the gauge bosons, the amplitude mode
with its closed form, the scan, the dynamics in time, Lord's
construction, every check of the Rust and the Python programs with its
deviation and tolerance, the data files with their SHA-256 checksums \\
\texttt{cover\_letter.pdf}, \texttt{.tex}, \texttt{.md} & the cover
letter to the editors \\
\texttt{submission\_info.pdf}, \texttt{.tex}, \texttt{.md} & what to
enter on the submission server: journal and article type, length, title
and byline, the abstract (typeset and as LaTeX source), the subject
classification (PhySH), the data availability statement, the comments
for the editor and the suggested referees, the files to upload, the
steps on the server, and the statements that only the author can
confirm \\
\texttt{submission/} & exactly the files to upload: the article as PDF
and as LaTeX source with its \texttt{.bbl} and its three figures in the
same folder, the Supplemental Material and the cover letter \\
\texttt{src/} & the templates of the four documents
(\texttt{*.template.tex}) and the bibliography \texttt{references.bib};
every number that a program computes is a placeholder
\texttt{\{\{NAME\}\}} in the templates \\
\texttt{scripts/build\_paper.py} & the program that fills the
placeholders and builds the four documents and the folder
\texttt{submission/} \\
\texttt{figures/} & the four figures of the article and of the
Supplemental Material, copied by the build from the folder
\texttt{figures/} of \texttt{Revision\_SM\_Higgs} \\
\texttt{generated/numbers.csv} & every placeholder with the value that
the last build inserted \\
\texttt{generated/build\_report.txt} & the checks of the last build
(Section 5) \\
\texttt{references.bib} & the copy of \texttt{src/references.bib} that
the build uses \\
\texttt{README.md}, \texttt{.tex}, \texttt{.pdf} & this description \\
\end{longtable}

\hypertarget{how-the-numbers-are-made}{%
\section{4. How the numbers are made}\label{how-the-numbers-are-made}}

The article prints results of the programs of the revision in the folder
\texttt{Revision\_SM\_Higgs}: the Rust programs \texttt{sh\_quantum},
\texttt{sh\_superfluid}, \texttt{sh\_higgs} and \texttt{sh\_dynamics},
which write the data files of \texttt{data/}; the Python programs
\texttt{python/sh\_quantum.py}, \texttt{sh\_superfluid.py},
\texttt{sh\_higgs.py} and \texttt{sh\_dynamics.py}, which check them
with their own engine and by other methods; and the Wolfram scripts
\texttt{wolfram/source/sh\_*.wls}, which check the group theory and the
closed forms in exact arithmetic, and whose checks Part R of the
notebook \texttt{superfluid.nb} repeats. Every check is written to a
report in \texttt{data/} with its deviation and its tolerance.

\texttt{scripts/build\_paper.py} takes the numbers that the manuscripts
of the revision state from the fragments that their fragment programs
wrote into \texttt{provenance/generated/} (\texttt{sh\_quantum.md},
\texttt{sh\_superfluid.md}, \texttt{sh\_higgs.md},
\texttt{sh\_dynamics.md}, \texttt{notebook.md}), so that the article
states exactly the numbers of the manuscripts; it computes the few
others from the data files (the cutoff in TeV, the ratio to the Standard
Model, the bound of \(\lvert\rho_{\rm SH}-1\rvert\), the masses of the
constituents in GeV). A bound of agreement is the largest deviation of
the checks it summarizes, rounded \textbf{up}, so that a stated bound is
never smaller than a measured deviation. The tables of results of the
Supplemental Material are the tables of the manuscripts, converted from
Markdown to LaTeX and set as floats, which never break across a page;
the tables of the checks are read from the reports. Where a statement of
a check cites the document of the hypotheses or an equation of a
manuscript of the stages, the build replaces the citation by the place
of the Supplemental Material that states the same thing (the
Supplemental Material refers to no other document and says so in its
Sec. S11), and it stops if a citation of another document remains. Typed
in the templates are only the parameters of the model and of the
numerical method, mathematical constants, the quantum numbers that the
hypotheses assign, and measured values with their references. The values
inserted by the last build are listed in \texttt{generated/numbers.csv}.

The build compiles the article and the Supplemental Material with
pdfLaTeX, BibTeX and REVTeX 4.2, the cover letter and the submission
information with pdfLaTeX, and writes the Markdown versions with pandoc.
It fixes the date that pdfTeX writes into a PDF to the date of the
documents (22 November 2026), so that a rebuild from the same sources
and data gives the same bytes.

To build (from the folder \texttt{Revision\_SM\_Higgs}, after the
programs have written \texttt{data/}, the fragment programs
\texttt{provenance/generated/}, and the figure programs
\texttt{python/sh\_article\_figures.py} and
\texttt{python/sh\_higgs\_figures.py} \texttt{figures/}):

\begin{Shaded}
\begin{Highlighting}[]
\BuiltInTok{cd} \StringTok{"/Users/nsh/Developer/github/spacetimes/Revision\_SM\_Higgs"}
\ExtensionTok{python3}\NormalTok{ python/sh\_article\_figures.py                 }\CommentTok{\# Figures 1 to 3 of the article}
\ExtensionTok{python3}\NormalTok{ python/sh\_higgs\_figures.py                   }\CommentTok{\# Figure S1 of the Supplemental Material}
\ExtensionTok{python3}\NormalTok{ paper\_prd/scripts/build\_paper.py             }\CommentTok{\# the four documents and submission/}
\FunctionTok{bash}\NormalTok{ code/scripts/build\_docs.sh paper\_prd/README.md  }\CommentTok{\# this README as .tex and .pdf}
\end{Highlighting}
\end{Shaded}

\texttt{build\_paper.py} needs Python 3, a TeX installation with REVTeX
4.2 (\texttt{revtex4-2.cls}, \texttt{apsrev4-2.bst}) and the LaTeX
packages \texttt{listings} and \texttt{longtable}, and pandoc. It runs
LaTeX again, up to three more times, while the labels of a document
still move. Its exit status is 1 if a document has an undefined or
multiply defined reference, labels that do not settle, a lost float or
an overfull box wider than 1 pt, or if the folder \texttt{submission/}
fails its check.

\hypertarget{the-checks-of-the-last-build}{%
\section{5. The checks of the last
build}\label{the-checks-of-the-last-build}}

\begin{verbatim}
build_paper.py: the checks of the builds
higgs_composite_qg: 7 pages; undefined: none; overfull boxes wider than 1 pt: none -> PASS
higgs_composite_qg_SM: 30 pages; undefined: none; overfull boxes wider than 1 pt: none -> PASS
cover_letter: 2 pages; undefined: none; overfull boxes wider than 1 pt: none -> PASS
submission/: the flat source compiles alone to 7 pages (the manuscript has 7); undefined: none -> PASS
submission_info: 5 pages; undefined: none; overfull boxes wider than 1 pt: none -> PASS
\end{verbatim}

The check of \texttt{submission/} compiles a copy of the folder alone in
a temporary folder, with the figures beside the source as the server
keeps them, and compares the number of pages with the article.

\hypertarget{corrections-made-before-this-version-was-kept}{%
\section{6. Corrections made before this version was
kept}\label{corrections-made-before-this-version-was-kept}}

Each of the following was found, by proofreading the built documents or
by a check of the build, before the first commit of this version of the
folder, and is corrected in it.

\begin{enumerate}
\def\labelenumi{\arabic{enumi}.}
\tightlist
\item
  \textbf{The masses of the constituents.} A draft said that the
  lightest fermionic excitations have masses of a few hundred GeV. The
  gap of the ground states of the second version is \(848\) to
  \(1.09\times10^{3}\) GeV, and the article states that range.
\item
  \textbf{The direction of the scan.} A draft said that the ratio
  \(M_H/v_{\rm SH}\) grows near the critical coupling. The scan shows
  the opposite: the ratio is smallest near the critical coupling and
  grows with the coupling, because the pair self-energy grows while the
  scale stays small. The build now checks that the ratio increases with
  the coupling at each mass of the second version, and it inserts the
  smallest ratio and the largest scale from the data.
\item
  \textbf{Statements wider than the computation.} Statements about the
  scan, the checks and the broken generators were first written for more
  cases than were computed; they now name the cases (the second version,
  the pairing-only methods, the generators that are symmetries of the
  version).
\item
  \textbf{The paths of the figures.} The source named the figures
  without their folder \texttt{figures/}; the check of the folder
  \texttt{submission/} found it. The source now names the folder, and
  the copy in \texttt{submission/} drops it, as the server keeps all
  files in one folder.
\item
  \textbf{Overfull lines.} The equation of the critical coupling was
  split, and the widths of the columns of the tables of the Supplemental
  Material were set.
\item
  \textbf{The legends.} The legends of Figures 1 and 2 covered curves
  and were moved.
\item
  \textbf{The notation of the tables.} The tables of the Supplemental
  Material first printed the symbols of the manuscripts (with the suffix
  SH where the article has none); the build now maps them to the symbols
  of the article, and the symbol table of the Supplemental Material
  gained the two potentials \(V^{(0)}\) and \(V^{(1)}\) that its tables
  use.
\item
  \textbf{An undefined mass.} The free Hamiltonian of the singlets in
  the Supplemental Material used a mass \(m\) that it did not define; it
  now names the masses of the two singlets.
\item
  \textbf{Where KS and CJT differ.} The Supplemental Material proves
  that KS and CJT coincide on the states of the paper, but its table of
  the normal vacuum shows the two methods differing in the response to
  the mixings of the fields. It now says why: such changes leave the
  span of the Kohn-Sham densities, and only CJT keeps their exchange
  contraction.
\item
  \textbf{Terms used before their definition.} The pairing-only methods
  (HS (0,1) and HS (0,0)) were named before they were defined; both
  documents now define them where the methods are introduced. The
  integral of the critical coupling was written \(I_1(m)\) for
  \(I_1(m_{\rm LL})\), and the stages that the tables of the checks name
  were not explained; Sec. S11 now explains them.
\item
  \textbf{The index of the hypercharge.} The mass matrix of the gauge
  bosons used the index 4 for the hypercharge while the eight generators
  of the global group were not listed; the eight generators are now
  named, and the hypercharge has the index \(Y\).
\item
  \textbf{The rank in the first version.} The Supplemental Material said
  that the matrix of the eight susceptibilities has rank three in the
  first version. There \(SU(2)_{\rm R}\) is not a symmetry, and the
  program counts the rank among \(T^1\), \(T^2\), \(T^3\),
  \(R_{\rm SH}\) and \(F^{\rm tot}_{\rm SH}\); the text now says so.
\item
  \textbf{A table under a wrong header.} The table of the scale broke
  across a page, and its last two rows appeared under the header of the
  next table, the table of the gap (a fault of long tables under REVTeX
  4.2 when a page breaks just before the next long table begins). The
  short tables are now floats, and the captions of the remaining long
  tables end their label with a period, as the floats do.
\item
  \textbf{Citations of other documents.} The statements of the checks
  cited the document of the hypotheses and equations of the manuscripts
  of the stages; the build replaces them (Section 4).
\item
  \textbf{The smallest ratio of the scan.} The abstract said that the
  ratio ``stays above'' its smallest value; it now says that the ratio
  is at least that value.
\item
  \textbf{The parameter \(\rho\) in the scan.} The article said that
  \(\rho_{\rm SH}=1\) at every point of the scan; that holds in the
  second version only (in the first version \(\rho_{\rm SH}\) grows with
  the coupling), and the range of the ratio that the sentence gives is
  that of the second version. The sentence now says so.
\item
  \textbf{The range of the factor.} The article said that the amplitude
  mode is 16 to 19 times too heavy without saying where the two ends
  come from; it now gives the factor of the linear response and that of
  the effective Lagrangian.
\item
  \textbf{Labels that move.} The build did not check that the labels had
  settled after its passes of LaTeX; it now runs more passes and fails
  if they do not settle.
\item
  \textbf{A plural.} The submission information wrote ``2 page(s)''; it
  now writes ``2 pages''.
\end{enumerate}

\hypertarget{provenance}{%
\section{7. Provenance}\label{provenance}}

The four documents, the folder \texttt{submission/},
\texttt{generated/numbers.csv} and \texttt{generated/build\_report.txt}
are written by \texttt{scripts/build\_paper.py} from the templates in
\texttt{src/}, the fragments of \texttt{provenance/generated/} and the
data files and reports of \texttt{data/}; the Supplemental Material
lists the data files and reports with their SHA-256 checksums. The
checksums of the templates, of the build program and of the files to
upload, at the time this README was built:

\begin{itemize}
\tightlist
\item
  \texttt{paper\_prd/scripts/build\_paper.py}:\\
  \texttt{cc7aa7a9b49f76ece537f8e8057aae936c8ee113d4bbe1e3bb966af55db65edd}
\item
  \texttt{paper\_prd/src/higgs\_composite\_qg.template.tex}:\\
  \texttt{98c3489ef1c8aec2bfb40392bbcbec625b1121c5d3bd86d6d9c264627ff3ba6e}
\item
  \texttt{paper\_prd/src/higgs\_composite\_qg\_SM.template.tex}:\\
  \texttt{9e2415a9c972794b2894ff2bdbd5eb94401332d46aa9b1af965fa78f5b3223a5}
\item
  \texttt{paper\_prd/src/cover\_letter.template.tex}:\\
  \texttt{def8931ca97b6d6cb7ffb150a4bfd8d311ede639becebf1f71f40f770884de56}
\item
  \texttt{paper\_prd/src/submission\_info.template.tex}:\\
  \texttt{8fbad56cae1684bdc8a64aad5e7f4e6f4cbff37500b78c380191c7dddea7402e}
\item
  \texttt{paper\_prd/src/references.bib}:\\
  \texttt{d0a061100956b7b41bf248b59210f9fadc6551aa9489884dd78f34c4d8a5c8fb}
\item
  \texttt{paper\_prd/submission/higgs\_composite\_qg.pdf}:\\
  \texttt{11f3480b122e922a54db4f6e82be2443ce675f713fbe8b7dc8e1f23fb22ce9e4}
\item
  \texttt{paper\_prd/submission/higgs\_composite\_qg.tex}:\\
  \texttt{9534b247f0d2f0ffdd48dcdfb4663a0aecfccef9ac818ed5a3dcb01b27758f62}
\item
  \texttt{paper\_prd/submission/higgs\_composite\_qg.bbl}:\\
  \texttt{68992b5d6923dd5ce72dc2eefc6204ba1c147e35f01cce11889567adc1f1ad82}
\item
  \texttt{paper\_prd/submission/higgs\_composite\_qg\_SM.pdf}:\\
  \texttt{8d0864eb25ce443f500dec24708e37cf652ec758bdd4f9b68de64544f2ed401d}
\item
  \texttt{paper\_prd/submission/cover\_letter.pdf}:\\
  \texttt{051e02feeb149593b317cc0b78ff916068b714c7e53346da2f53c72b17c81a7c}
\end{itemize}

\end{document}
